\documentclass{article}
\usepackage{amsmath,amssymb,amsthm,mathtools}
\newtheorem{theorem}{Theorem}
\newtheorem{lemma}{Lemma}
\newtheorem{proposition}{Proposition}
\newtheorem{corollary}{Corollary}
\newtheorem{definition}{Definition}
\newtheorem{example}{Example}
\title{ShadowBench algebra/L2/alg\_grob\_L2\_013}
\begin{document}
\maketitle
% Problem id: algebra/L2/alg_grob_L2_013
% Companion instructions: docs/instructions.md
% Candidate skeletons: docs/skeletons/
% Lean target: ShadowBench/Source/Main.lean

Lemma (lm_add_le_of_both_lm_le_mon)
Upper bound of the leading monomial of addition

Let $f_1, f_2 \in R[\sigma]$ be multivariate polynomials (which might be zero). Here, $\sigma$ is a set of variable symbols over which a monomial order $>$ is fixed, and $R$ is a nontrivial commutative semiring. Given a monomial $x^\delta$ such that $\mathrm{LM}(f_1) \le x^\delta$ and $\mathrm{LM}(f_2) \le x^\delta$, we have $\mathrm{LM}(f_1 + f_2) \le x^\delta$.

Proof)
It suffices to take a polynomial $g \in R[\sigma]$ whose leading monomial is $\delta$. Such polynomial exists from $R \ne \{0\}$; in particular, $g$ can be taken as the monomial $cx^\delta$ for some $c \in R \setminus \{0\}$. Now we obtain the result from the same inequality for $f_1, f_2$ and $g$.
\end{document}
